\section{Conclusiones principales}

El trabajo desarrollado permite concluir que es posible construir una plataforma
de visualizaci\'on y apoyo al an\'alisis de estrategias de trading. Esta plataforma
combina formalizaci\'on de reglas, evaluaci\'on hist\'orica, consulta de datos,
seguimiento del estado del sistema y demostraci\'on controlada con MT5. La
aportaci\'on principal no reside en una \'unica estrategia ni en un resultado
aislado. Reside en el flujo completo: partir de datos de mercado, convertir
criterios t\'ecnicos en reglas calculables, compararlas frente a referencias,
organizar las evidencias generadas y presentarlas en una interfaz que facilite
la revisi\'on supervisada.

En la parte cuantitativa, la estrategia \texttt{macd\_breakout} es la que ofrece
la lectura m\'as defendible dentro del benchmark final de comparaci\'on. Sus
resultados agregados son m\'as favorables que los de las estrategias de
referencia evaluadas bajo el mismo marco, aunque tambi\'en muestran variabilidad
entre grupos de activos y timeframes. Por ello, la conclusi\'on no debe
formularse como una garant\'ia de rentabilidad futura, sino como una evidencia de
mejor comportamiento relativo en los experimentos realizados. La estrategia
\texttt{fib\_limit}, en cambio, queda como una formalizaci\'on m\'as irregular:
permite estudiar una idea t\'ecnica concreta, pero no alcanza una lectura
agregada suficientemente estable.

La l\'inea basada en el enfoque de Men\'endez cumple una funci\'on distinta. Sus
resultados no son positivos desde el punto de vista operativo, pero s\'i aportan
valor metodol\'ogico al mostrar las dificultades de traducir una lectura visual
de Elliott a reglas cuantificables. Esta parte del trabajo ayuda a separar la
inspiraci\'on t\'ecnica de la validaci\'on experimental. Una teor\'ia puede ser
\'util para estructurar la revisi\'on del mercado, pero no debe tratarse como
se\~nal operativa sin una prueba separada y reproducible.

WaveCount y WeaveCount refuerzan esa misma separaci\'on. Su valor est\'a en
convertir una lectura estructural del precio en informaci\'on revisable, no en
generar \'ordenes. Las etiquetas de onda, los niveles de calidad visual y las
advertencias de estudio permiten priorizar casos y entender mejor el contexto,
pero se mantienen fuera de la ejecuci\'on. Este matiz es importante porque
sit\'ua el an\'alisis estructural como una capa de apoyo, no como un mecanismo de
decisi\'on autom\'atica.

La plataforma final integra estos resultados en un entorno de visualizaci\'on y
trazabilidad m\'as ordenado. El Trading Center Dashboard, el Screener, el AI Analyst, MT5 Bot,
RiskGuard y Telegram no tienen todos el mismo papel, pero juntos permiten
reducir la dispersi\'on de informaci\'on y mejorar la trazabilidad del proceso. El
dashboard centraliza la revisi\'on, el Screener prioriza candidatos de estudio, el
AI Analyst prepara informes, RiskGuard audita condiciones de riesgo, MT5 permite
lectura y ciclo demo controlado, y Telegram queda como canal informativo. La
separaci\'on entre estas capas es uno de los resultados m\'as relevantes del
trabajo, porque evita confundir revisi\'on, contexto, validaci\'on, orden demo y
operaci\'on real.

\section{Cumplimiento de los objetivos}

El objetivo general se considera alcanzado dentro del alcance definido. Se ha
construido una base t\'ecnica que permite estudiar estrategias de trading,
ordenar evidencias y presentar informaci\'on relevante para una toma de
decisiones supervisada. El sistema no sustituye el criterio del usuario ni
pretende operar de forma totalmente aut\'onoma. S\'i proporciona un flujo de
trabajo que reduce la carga de revisi\'on inicial y ayuda a aprovechar mejor los
datos disponibles. Tambi\'en permite relacionar cada conclusi\'on
con datos, reglas, tablas, gr\'aficos o registros generados durante el proceso.

Respecto a la preparaci\'on de datos, se ha desarrollado un flujo de ingesta y
actualizaci\'on orientado a trabajar con series intradiarias y diarias. La
informaci\'on externa de contexto se mantiene separada del conjunto principal de
mercado, especialmente cuando depende de disponibilidad, coste o forma de
acceso. Esta decisi\'on permite conservar una base reproducible para los
experimentos y, al mismo tiempo, incorporar contexto adicional cuando resulta
necesario para la revisi\'on.

El objetivo de formalizar estrategias tambi\'en queda cubierto. Las reglas
basadas en \texttt{macd\_breakout}, \texttt{fib\_limit} y RSI se han traducido a
condiciones calculables, diferenciando entre setup, contexto y se\~nal de
estudio. Adem\'as, se ha evaluado una l\'inea espec\'ifica basada en el enfoque de
Men\'endez, lo que permite documentar tanto los aciertos como las limitaciones de
formalizar una lectura t\'ecnica m\'as discrecional.

En relaci\'on con las pruebas hist\'oricas, el trabajo ha permitido comparar las
estrategias formalizadas frente a referencias sencillas y analizar sus
resultados por bloques independientes. Esta evaluaci\'on se ha planteado con
prudencia, evitando interpretar las tablas como una cartera \'unica o como una
promesa de rentabilidad futura. Tambi\'en se ha prestado atenci\'on al control de
\textit{look-ahead}, especialmente en las reglas donde la detecci\'on de
estructura o niveles t\'ecnicos pod\'ia introducir informaci\'on futura si no se
trataba con cuidado.

El objetivo de incorporar an\'alisis estructural se ha materializado mediante
WaveCount y WeaveCount. Estas herramientas no cierran una se\~nal operativa,
pero s\'i ofrecen una forma revisable de representar hip\'otesis de onda, calidad
visual y contexto. De esta forma, la estructura del precio se incorpora al
sistema sin mezclarse con la ejecuci\'on de \'ordenes.

La plataforma interactiva queda implementada como herramienta de visualizaci\'on y revisi\'on. El
dashboard permite consultar mercado, correlaciones, WeaveCount, Screener, MT5
Bot y AI Analyst desde una interfaz com\'un. Adem\'as, mantiene una separaci\'on
clara: observa, resume y organiza informaci\'on, pero no ejecuta operaciones
desde el panel.

Por \'ultimo, el objetivo relacionado con MT5, RiskGuard y Telegram se cumple
como demostraci\'on controlada. Se ha comprobado un flujo demo con preparaci\'on de
orden, env\'io auditado, lectura de posici\'on, cierre y comprobaci\'on final.
RiskGuard queda como capa de control previo y Telegram como canal informativo de
observabilidad. Esta integraci\'on no habilita operaci\'on real ni convierte el
sistema en una soluci\'on preparada para ejecuci\'on sin supervisi\'on.

\section{Limitaciones}

La primera limitaci\'on est\'a asociada a la interpretaci\'on de los resultados
hist\'oricos. Aunque el benchmark permite comparar estrategias bajo un marco
com\'un, los resultados dependen de los activos, timeframes, reglas de
entrada, reglas de salida y supuestos de ejecuci\'on considerados. Por tanto, no
pueden extrapolarse directamente a rentabilidad futura ni a condiciones reales
de mercado sin validaciones adicionales.

Una segunda limitaci\'on es que el sistema no se ha evaluado como servicio
continuo durante un periodo prolongado. La arquitectura permite actualizar
archivos generados y alimentar el dashboard en entorno local. Para
considerarla una soluci\'on operativa estable har\'ian falta pruebas de
disponibilidad, monitorizaci\'on, recuperaci\'on ante errores, control de fuentes
externas y despliegue del planificador como servicio mantenido.

La tercera limitaci\'on afecta a la parte operativa. MT5 Bot y RiskGuard se han
cerrado como ciclo demo controlado, no como bot de trading real. La ejecuci\'on en
cuenta real, la autonom\'ia completa, la gesti\'on integral del riesgo bajo
condiciones reales y la gesti\'on de incidencias quedan fuera del alcance. La
capa de riesgo permite auditar \'ordenes demo y controlar condiciones b\'asicas de
exposici\'on, pero no sustituye a una gesti\'on de cartera que considere
diversificaci\'on, correlaciones, concentraci\'on por activo o divisa y dependencia
entre mercados. Del mismo modo, Telegram informa sobre eventos y estados, pero
no funciona como consola, no confirma operaciones y no env\'ia \'ordenes.

La cuarta limitaci\'on se relaciona con el AI Analyst. El asistente puede ayudar a
preparar informes y sintetizar informaci\'on, pero sus salidas deben entenderse
como apoyo a la revisi\'on. No sustituye a RiskGuard, no aprueba operaciones y no
convierte una se\~nal de estudio en una decisi\'on de trading. Esta separaci\'on es
necesaria para mantener el sistema dentro de un uso supervisado y defendible.

Finalmente, la informaci\'on macroecon\'omica y externa de contexto no queda
integrada como un conjunto de datos estructurado equivalente a las series de
mercado. En el trabajo se utiliza como apoyo a la revisi\'on cuando procede,
especialmente a trav\'es del AI Analyst, pero no como una fuente sistem\'atica
integrada en todos los experimentos. Esta decisi\'on evita problemas de
disponibilidad y acceso, aunque limita la profundidad con la que puede
estudiarse el impacto de eventos externos.

\section{Trabajo futuro}

Una primera l\'inea futura consiste en ampliar la validaci\'on estad\'istica de las
estrategias formalizadas. Ser\'ia conveniente estudiar nuevos periodos, m\'as
activos, validaciones sobre periodos no utilizados para definir las reglas,
pruebas de sensibilidad y an\'alisis de costes m\'as detallados. Esto permitir\'ia
comprobar si el comportamiento relativo de \texttt{macd\_breakout} se mantiene
fuera del corte analizado y si \texttt{fib\_limit} puede mejorar mediante
condiciones de contexto m\'as selectivas.

Otra l\'inea relevante es incorporar nuevas estrategias al Screener siguiendo la
misma separaci\'on entre regla principal, contexto y se\~nal de estudio. El
sistema se ha dise\~nado como una base ampliable, capaz de crecer con nuevas
estrategias, m\'as periodos de prueba y estad\'isticas acumuladas. Esta
escalabilidad no convierte el sistema en una soluci\'on operativa final, pero s\'i
lo deja preparado para evolucionar desde reglas aisladas hacia un proceso de
an\'alisis m\'as ordenado y verificable. En ese crecimiento, el l\'imite de
seguridad debe mantenerse: primero se formaliza, despu\'es se eval\'ua y solo m\'as
adelante se estudia si tiene sentido plantear una demostraci\'on operativa.

Una extensi\'on m\'as ambiciosa consistir\'ia en estudiar modelos de aprendizaje
autom\'atico o redes neuronales como comparativa metodol\'ogica. En ese caso, el
reto no ser\'ia solo entrenar un modelo, sino definir etiquetas fiables, separar
adecuadamente entrenamiento y evaluaci\'on, controlar el sobreajuste y justificar
qu\'e informaci\'on estar\'ia disponible en el momento de la predicci\'on. Por ello,
esta l\'inea se deja como trabajo futuro y no como parte de la validaci\'on
principal de esta memoria.

Tambi\'en queda margen para profundizar en WaveCount y WeaveCount. Como trabajo
futuro, podr\'ian estudiarse formas de medir la utilidad del contexto estructural
en la revisi\'on de setups, sin convertirlo necesariamente en filtro autom\'atico.
Esto permitir\'ia analizar si ciertas configuraciones de onda ayudan a priorizar
mejor los casos o a descartar situaciones ambiguas.

Desde el punto de vista de la plataforma, el siguiente paso ser\'ia comprobar el
sistema durante periodos m\'as largos y en un despliegue m\'as parecido a un uso
continuado. Para ello ser\'ia necesario desplegar el planificador como servicio,
monitorizar errores de ingesta, controlar tiempos de actualizaci\'on y generar alertas
de fallos. Tambi\'en habr\'ia que mantener una selecci\'on clara de tablas, capturas
e informes que permitan reproducir y revisar el estado del sistema.

En la parte operativa, una evoluci\'on razonable pasar\'ia por reforzar RiskGuard,
probar durante m\'as tiempo el ciclo demo y revisar la gesti\'on de posiciones
simult\'aneas. Tambi\'en convendr\'ia estudiar m\'etricas de exposici\'on en condiciones m\'as cercanas a
un uso real. Ser\'ia necesario profundizar adem\'as en la gesti\'on de cartera:
correlaciones, concentraci\'on por divisa o sector, modelos de dependencia entre
activos y relaci\'on con eventos macroecon\'omicos. Cualquier avance hacia
operaci\'on real deber\'ia mantenerse separado, documentado y sujeto a controles
adicionales.

Por \'ultimo, el AI Analyst y Telegram podr\'ian ampliarse como herramientas de
apoyo al seguimiento. El AI Analyst podr\'ia generar informes comparables entre
setups y sesiones, mientras que Telegram podr\'ia enviar res\'umenes informativos
m\'as completos sobre estado del sistema, posiciones demo, errores y revisiones
disponibles. En ambos casos, la evoluci\'on deber\'ia conservar el l\'imite
definido en este trabajo: informar y ayudar a revisar, no decidir ni ejecutar.
